\chapter{2015年北京卷（文）}

\section{单选题}

\begin{problem}
若集合 \(A = \{x \mid -5 < x < 2\}\)，\(B = \{x \mid -3 < x < 3\}\)，则 \(A \cap B =\)（\quad）

\choices
    {\(\{x\mid -3 < x < 2\}\)}
    {\(\{x\mid -5 < x < 2\}\)}
    {\(\{x\mid -3 < x < 3\}\)}
    {\(\{x\mid -5 < x < 3\}\)}
\end{problem}

\begin{answer}
A
\end{answer}

\begin{solution}
同时满足 \(-5<x<2\) 和 \(-3<x<3\) 的实数，必须满足下界取两者中较大的 \(-3\)，上界取两者中较小的 \(2\)，所以 \(A\cap B=\{x\mid -3<x<2\}\)，故选 A.
\end{solution}

\begin{problem}
圆心为 \((1,1)\) 且过原点的圆的方程是（\quad）

\choices
    {\((x - 1)^2 + (y - 1)^2 = 1\)}
    {\((x + 1)^{2} + (y + 1)^{2} = 1\)}
    {\((x + 1)^{2} + (y + 1)^{2} = 2\)}
    {\((x - 1)^{2} + (y - 1)^{2} = 2\)}
\end{problem}

\begin{answer}
D
\end{answer}

\begin{solution}
圆心为 \((1,1)\)，且经过原点，因此半径的平方为 \((1-0)^2+(1-0)^2=2\). 所以圆的方程为 \((x-1)^2+(y-1)^2=2\)，故选 D.
\end{solution}

\begin{problem}
下列函数中为偶函数的是（\quad）

\choices
    {\(y = x^{2}\sin x\)}
    {\( y = x^{2}\cos x \)}
    {\(y = |\ln x|\)}
    {\( y = 2^{-x} \)}
\end{problem}

\begin{answer}
B
\end{answer}

\begin{solution}
选项 A 的函数定义域为 \(\R\)，且 \(f(-x)=(-x)^2\sin(-x)=-x^2\sin x=-f(x)\)，是奇函数；选项 B 的函数定义域为 \(\R\)，且 \(f(-x)=(-x)^2\cos(-x)=x^2\cos x=f(x)\)，是偶函数. 选项 C 的定义域为 \((0,+\infty)\)，不是关于原点对称的集合；选项 D 满足 \(f(-x)=2^x\ne 2^{-x}=f(x)\)（例如取 \(x=1\)）. 故选 B.
\end{solution}

\begin{problem}
某校老年，中年和青年教师的人数见下表，采用分层抽样的方法调查教师的身体状况，在抽取的样本中，青年教师有\(320\text{ 人}\)，则该样本中的老年教师人数为（\quad）

\begin{center}
\begin{tabular}{|c|c|}
\hline
类别 & 人数 \\
\hline
老年教师 & \(900\) \\
\hline
中年教师 & \(1800\) \\
\hline
青年教师 & \(1600\) \\
\hline
合计 & \(4300\) \\
\hline
\end{tabular}
\end{center}

\choices
    {\(90\)}
    {\(100\)}
    {\(180\)}
    {\(300\)}
\end{problem}

\begin{answer}
C
\end{answer}

\begin{solution}
分层抽样要求各层抽取比例相同. 青年教师总体有 \(1600\) 人，样本中有 \(320\) 人，所以抽样比例为 \(320/1600=1/5\). 因此样本中的老年教师人数为 \(900\times\frac15=180\)，故选 C.
\end{solution}

\begin{problem}
执行如图所示的程序框图，输出 \(k\) 的值为（\quad）

\bitmapfigure[max width=0.48\linewidth]{img/2015/beijing_liberal/q5_fig1.png}

\choices
    {\(3\)}
    {\(4\)}
    {\(5\)}
    {\(6\)}
\end{problem}

\begin{answer}
B
\end{answer}

\begin{solution}
程序开始时 \(k=0\)，\(a=3\)，\(q=\frac12\). 每循环一次先令 \(a\leftarrow aq\)，再令 \(k\leftarrow k+1\)，然后判断 \(a<\frac14\). 四次循环后的数值依次为 \(\left(a,k\right)=\left(\frac32,1\right)\)、\(\left(\frac34,2\right)\)、\(\left(\frac38,3\right)\)、\(\left(\frac{3}{16},4\right)\). 前三次均不满足 \(a<\frac14\)，第四次满足 \(\frac{3}{16}<\frac14\)，所以输出 \(k=4\)，故选 B.
\end{solution}

\begin{problem}
设 \(\bs{a}\)，\(\bs{b}\) 是非零向量，“\(\bs{a} \cdot \bs{b} = |\bs{a}| |\bs{b}|\)”是“\(\bs{a} \parallel \bs{b}\)”的（\quad）

\choices
    {充分而不必要条件}
    {必要而不充分条件}
    {充分必要条件}
    {既不充分也不必要条件}
\end{problem}

\begin{answer}
A
\end{answer}

\begin{solution}
设非零向量 \(\bs{a}\) 与 \(\bs{b}\) 的夹角为 \(\theta\)，则 \(0\leq\theta\leq\pi\)，且 \(\bs{a}\cdot\bs{b}=|\bs{a}|\,|\bs{b}|\cos\theta\). 若 \(\bs{a}\cdot\bs{b}=|\bs{a}|\,|\bs{b}|\)，由于两向量均非零，可得 \(\cos\theta=1\)，从而 \(\theta=0\)，所以 \(\bs{a}\parallel\bs{b}\). 反过来，若 \(\bs{a}\parallel\bs{b}\)，两向量也可能方向相反，此时 \(\theta=\pi\)，有 \(\bs{a}\cdot\bs{b}=-|\bs{a}|\,|\bs{b}|\)，不满足题设等式. 因此题设条件是 \(\bs{a}\parallel\bs{b}\) 的充分而不必要条件，故选 A.
\end{solution}

\begin{problem}
某四棱锥的三视图如图所示，该四棱锥最长棱的棱长为（\quad）

\bitmapfigure[max width=0.62\linewidth]{img/2015/beijing_liberal/q7_fig1.png}

\choices
    {\(1\)}
    {\(\sqrt{2}\)}
    {\(\sqrt{3}\)}
    {\(2\)}
\end{problem}

\begin{answer}
C
\end{answer}

\begin{solution}
由三视图可设底面为边长为 \(1\) 的正方形，四个顶点记为 \(A\)，\(B\)，\(C\)，\(D\)，棱锥顶点为 \(S\). 三视图表明 \(SC\perp\) 底面且 \(SC=1\)，因此底面对角线 \(CA=\sqrt{2}\). 于是 \(SA^2=SC^2+CA^2=1+2=3\)，即 \(SA=\sqrt{3}\). 另外，顶点到底面相邻顶点的棱长为 \(\sqrt{1^2+1^2}=\sqrt{2}\)，底面棱长为 \(1\)，所以最长棱为 \(SA=\sqrt{3}\)，故选 C.
\end{solution}

\begin{problem}
某辆汽车每次加油都把油箱加满，下表记录了该车相邻两次加油时的情况.

\begin{center}
\begin{tabular}{|c|c|c|}
\hline
加油时间 & 加油量（升） & 加油时的累计里程（千米） \\
\hline
2015年5月1日 & \(12\) & \(35000\) \\
\hline
2015年5月15日 & \(48\) & \(35600\) \\
\hline
\end{tabular}
\end{center}

注：“累计里程”指汽车从出厂开始累计行驶的路程.

在这段时间内，该车每\(100\text{ 千米}\)平均耗油量为（\quad）

\choices
    {\(6\text{ 升}\)}
    {\(8\text{ 升}\)}
    {\(10\text{ 升}\)}
    {\(12\text{ 升}\)}
\end{problem}

\begin{answer}
B
\end{answer}

\begin{solution}
5 月 1 日加油后油箱已加满，5 月 15 日再次加满时加入的 \(48\text{ 升}\) 正好是这段时间内消耗的汽油量. 这段时间行驶的路程为 \(35600-35000=600\text{ 千米}\)，所以每 \(100\text{ 千米}\)的平均耗油量为 \(48\times\frac{100}{600}=8\text{ 升}\)，故选 B.
\end{solution}

\section{填空题}

\begin{problem}
复数 \(\i(1+\i)\) 的实部为\fillinblank{}.
\end{problem}

\begin{answer}
\(-1\)
\end{answer}

\begin{solution}
\(\i(1+\i)=\i+\i^2=\i-1=-1+\i\)，所以其实部为 \(-1\).
\end{solution}

\begin{problem}
\(2^{-3}\)，\(3^{\frac{1}{2}}\)，\(\log_2 5\) 三个数中最大的数是\fillinblank{}.
\end{problem}

\begin{answer}
\(\log_2 5\)
\end{answer}

\begin{solution}
\(2^{-3}=\frac18\)，所以 \(0<2^{-3}<1\). 又因为 \(1<3<4\)，所以 \(1<\sqrt{3}<2\). 由于底数 \(2>1\)，对数函数单调递增，且 \(5>4\)，故 \(\log_2 5>\log_2 4=2\). 因而 \(\log_2 5\) 最大.
\end{solution}

\begin{problem}
在 \(\triangle ABC\) 中，\(a = 3\)，\(b = \sqrt{6}\)，\(\angle A = \frac{2\pi}{3}\)，则 \(\angle B =\fillinblank{}\).
\end{problem}

\begin{answer}
\(\dfrac{\pi}{4}\)
\end{answer}

\begin{solution}
由正弦定理，\(\dfrac{a}{\sin A}=\dfrac{b}{\sin B}\)，所以
\(\sin B=\dfrac{b\sin A}{a}=\dfrac{\sqrt{6}\sin\frac{2\pi}{3}}{3}=\dfrac{\sqrt{6}\cdot\frac{\sqrt{3}}2}{3}=\frac{\sqrt{2}}2\). 又因为 \(A+B<\pi\)，所以 \(0<B<\pi-A=\frac{\pi}{3}\). 在此范围内满足 \(\sin B=\frac{\sqrt2}{2}\) 的角只有 \(B=\frac{\pi}{4}\).
\end{solution}

\begin{problem}
已知 \((2,0)\) 是双曲线 \(x^{2} - \frac{y^{2}}{b^{2}} = 1\)（\(b > 0\)）的一个焦点，则 \(b =\)\fillinblank{}.
\end{problem}

\begin{answer}
\(\sqrt{3}\)
\end{answer}

\begin{solution}
双曲线可写为 \(\dfrac{x^2}{1}-\dfrac{y^2}{b^2}=1\)，故其半实轴为 \(a=1\)，焦半距满足 \(c^2=a^2+b^2=1+b^2\). 已知焦点为 \((2,0)\)，所以 \(c=2\)，从而 \(1+b^2=4\). 由 \(b>0\)，得 \(b=\sqrt3\).
\end{solution}

\begin{problem}
如图，\(\triangle ABC\) 及其内部的点组成的集合记为 \(D\)，\(P(x, y)\) 为 \(D\) 中任意一点，则 \(z = 2x + 3y\) 的最大值为\fillinblank{}.

\bitmapfigure[max width=0.45\linewidth]{img/2015/beijing_liberal/q13_fig1.png}
\end{problem}

\begin{answer}
\(7\)
\end{answer}

\begin{solution}
由图可读出三角形的三个顶点为 \(A(2,1)\)、\(B(1,0)\)、\(C(0,2)\). 线性函数 \(z=2x+3y\) 在三角形及其内部的最大值出现在顶点处，分别计算得
\(z(A)=2\times2+3\times1=7\)，\(z(B)=2\times1+3\times0=2\)，\(z(C)=2\times0+3\times2=6\). 因此 \(z\) 的最大值为 \(7\).
\end{solution}

\begin{problem}
高三年级\(267\text{ 位}\)学生参加期末考试，某班\(37\text{ 位}\)学生的语文成绩，数学成绩与总成绩在全年级中的排名情况如图所示，甲，乙，丙为该班三位学生.

从这次考试成绩看，

\begin{circlelist}
\item 在甲，乙两人中，其语文成绩名次比其总成绩名次靠前的学生是\fillinblank{}；
\item 在语文和数学两个科目中，丙同学的成绩名次更靠前的科目是\fillinblank{}.
\end{circlelist}

\bitmapfigure[width=0.84\linewidth]{img/2015/beijing_liberal/q14_fig1.png}
\end{problem}

\begin{answer}
乙；数学
\end{answer}

\begin{solution}
两个散点图的横坐标都是总成绩名次，纵坐标分别是语文成绩名次和数学成绩名次，且名次数字越小表示名次越靠前. 在左图中比较甲、乙两点的横、纵坐标可知，乙点的纵坐标小于横坐标，而甲点的纵坐标大于横坐标，因此乙的语文成绩名次比其总成绩名次靠前. 丙在两个图中的总成绩名次相同；在左图中找到与右图丙点横坐标相同的对应点，可见其语文名次纵坐标大于右图丙点的数学名次纵坐标，所以丙的数学成绩名次更靠前. 故两空依次为乙、数学.
\end{solution}

\section{解答题}

\begin{problem}
已知函数 \( f(x) = \sin x - 2\sqrt{3}\sin^2\frac{x}{2} \).

\begin{enumerate}
\item 求 \( f(x) \) 的最小正周期；
\item 求 \( f(x) \) 在区间 \(\left[0, \frac{2\pi}{3}\right]\) 上的最小值.
\end{enumerate}
\end{problem}

\begin{solution}
\begin{enumerate}
\item 由半角公式 \(\sin^2\frac{x}{2}=\frac{1-\cos x}{2}\)，得
\[
f(x)=\sin x-\sqrt{3}(1-\cos x)=\sin x+\sqrt{3}\cos x-\sqrt{3}=2\sin\left(x+\frac{\pi}{3}\right)-\sqrt{3}
\]
正弦函数的最小正周期为 \(2\pi\)，平移自变量和加减常数都不改变周期，因此 \(f(x)\) 的最小正周期为 \(2\pi\).
\item 当 \(x\in\left[0,\frac{2\pi}{3}\right]\) 时，\(x+\frac{\pi}{3}\in\left[\frac{\pi}{3},\pi\right]\). 在区间 \(\left[\frac{\pi}{3},\pi\right]\) 上，\(\sin\left(x+\frac{\pi}{3}\right)\) 的最小值为 \(0\)，在 \(x=\frac{2\pi}{3}\) 时取得. 所以
\[
f(x)=2\sin\left(x+\frac{\pi}{3}\right)-\sqrt{3}\geq-\sqrt{3}
\]
且等号可以取得，故 \(f(x)\) 在该区间上的最小值为 \(-\sqrt{3}\).
\end{enumerate}
\end{solution}

\begin{problem}
已知等差数列 \(\{a_{n}\}\) 满足 \(a_1 + a_2 = 10\)，\(a_4 - a_3 = 2\).

\begin{enumerate}
\item 求 \(\{a_{n}\}\) 的通项公式；
\item 设等比数列 \(\{b_{n}\}\) 满足 \(b_{2}=a_{3}\)，\(b_{3}=a_{7}\)，问： \(b_{6}\) 与数列 \(\{a_{n}\}\) 的第几项相等？
\end{enumerate}
\end{problem}

\begin{solution}
\begin{enumerate}
\item 设等差数列 \(\{a_n\}\) 的公差为 \(d\). 由 \(a_4-a_3=d=2\)，以及 \(a_1+a_2=a_1+(a_1+d)=10\)，得 \(2a_1+2=10\)，所以 \(a_1=4\). 因此
\(a_n=a_1+(n-1)d=4+2(n-1)=2n+2\)，\((n=1,2,\ldots)\).
\item 由上式得 \(a_3=8\)，\(a_7=16\)，所以 \(b_2=8\)，\(b_3=16\). 设等比数列 \(\{b_n\}\) 的公比为 \(q\)，则 \(q=b_3/b_2=2\)，进而 \(b_1=b_2/q=4\). 因此 \(b_6=b_1q^5=4\cdot2^5=128\). 令 \(b_6=a_n\)，得 \(2n+2=128\)，解得 \(n=63\). 所以 \(b_6\) 与数列 \(\{a_n\}\) 的第 \(63\) 项相等.
\end{enumerate}
\end{solution}

\begin{problem}
某超市随机选取\(1000\text{ 位}\)顾客，记录了他们购买甲，乙，丙，丁四种商品的情况，整理成如下统计表，其中“√”表示购买，“×”表示未购买.

\begin{center}
\renewcommand{\arraystretch}{1.2}
\begin{tabular}{c|cccc}
\hline
顾客人数 & \multicolumn{4}{c}{商品} \\
\cline{2-5}
 & 甲 & 乙 & 丙 & 丁 \\
\hline
\(100\) & √ & × & √ & √ \\
\(217\) & × & √ & × & √ \\
\(200\) & √ & √ & √ & × \\
\(300\) & √ & × & √ & × \\
\(85\) & √ & × & × & × \\
\(98\) & × & √ & × & × \\
\hline
\end{tabular}
\end{center}

\begin{enumerate}
\item 估计顾客同时购买乙和丙的概率；
\item 估计顾客在甲，乙，丙，丁中同时购买\(3\text{ 种}\)商品的概率；
\item 如果顾客购买了甲，则该顾客同时购买乙，丙，丁中哪种商品的可能性最大？
\end{enumerate}
\end{problem}

\begin{solution}
\begin{enumerate}
\item 同时购买乙和丙的顾客对应表中同时标有乙、丙购买符号的行，只有人数为 \(200\) 的一行. 因此估计该概率为 \(\frac{200}{1000}=0.2\).
\item 同时购买 \(3\) 种商品的顾客对应人数为 \(100\) 和 \(200\) 的两行，共有 \(100+200=300\) 人. 因此估计该概率为 \(\frac{300}{1000}=0.3\).
\item 购买甲的顾客人数为 \(100+200+300+85=685\). 在这些顾客中，同时购买乙的有 \(200\) 人，同时购买丙的有 \(100+200+300=600\) 人，同时购买丁的有 \(100\) 人. 相应的条件概率分别为
\(\frac{200}{685}\)，\(\frac{600}{685}\)，\(\frac{100}{685}\).
其中 \(\frac{600}{685}\) 最大，所以可能性最大的是丙.
\end{enumerate}
\end{solution}

\begin{problem}
如图，在三棱锥 \(V-ABC\) 中，平面 \(VAB \perp\) 平面 \(ABC\)，\(\triangle VAB\) 为等边三角形，\(AC \perp BC\) 且 \(AC = BC = \sqrt{2}\)，\(O\)，\(M\) 分别为 \(AB\)，\(VA\) 的中点.

\begin{enumerate}
\item 求证：\(VB\parallel\) 平面 \(MOC\)；
\item 求证：平面 \(MOC\perp\) 平面 \(VAB\)；
\item 求三棱锥 \(V-ABC\) 的体积.
\end{enumerate}

\bitmapfigure[max width=0.62\linewidth]{img/2015/beijing_liberal/q18_fig1.png}
\end{problem}

\begin{solution}
\begin{enumerate}
\item 在三角形 \(VAB\) 中，\(O\)、\(M\) 分别是 \(AB\)、\(VA\) 的中点，所以由三角形中位线定理得 \(OM\parallel VB\). 又因为 \(M\notin\) 平面 \(ABC\)，平面 \(MOC\) 与平面 \(ABC\) 的交线为 \(OC\)，而 \(B\notin OC\)，所以 \(VB\) 不在平面 \(MOC\) 内. 直线 \(VB\) 平行于平面 \(MOC\) 内的直线 \(OM\)，且不在该平面内，故 \(VB\parallel\) 平面 \(MOC\).
\item 因为 \(AC=BC\)，且 \(O\) 是 \(AB\) 的中点，所以 \(C\) 和 \(O\) 都在 \(AB\) 的垂直平分线上，从而 \(OC\perp AB\). 平面 \(VAB\) 与平面 \(ABC\) 的交线为 \(AB\)，且两平面垂直；因此平面 \(ABC\) 内垂直于 \(AB\) 的直线 \(OC\) 垂直于平面 \(VAB\). 由于 \(OC\subset\) 平面 \(MOC\)，所以平面 \(MOC\perp\) 平面 \(VAB\).
\item 在直角三角形 \(ABC\) 中，\(AC=BC=\sqrt2\)，所以 \(AB=\sqrt{AC^2+BC^2}=2\)，且
\[
S_{\triangle ABC}=\frac12 AC\cdot BC=\frac12\cdot\sqrt2\cdot\sqrt2=1
\]
三角形 \(VAB\) 为边长 \(2\) 的等边三角形，\(O\) 为 \(AB\) 的中点，所以 \(VO\perp AB\)，且
\[
VO=\sqrt{2^2-\left(\frac{2}{2}\right)^2}=\sqrt3
\]
由于平面 \(VAB\perp\) 平面 \(ABC\)，且 \(VO\) 在平面 \(VAB\) 内垂直于交线 \(AB\)，故 \(VO\perp\) 平面 \(ABC\)，它就是三棱锥 \(V-ABC\) 的高. 因此
\[
V_{V-ABC}=\frac13S_{\triangle ABC}\cdot VO=\frac13\cdot1\cdot\sqrt3=\frac{\sqrt3}{3}
\]
\end{enumerate}
\end{solution}

\begin{problem}
设函数 \(f(x) = \frac{x^2}{2} - k\ln x\)，\(k > 0\).

\begin{enumerate}
\item 求 \( f(x) \) 的单调区间和极值；
\item 证明：若 \( f(x) \) 存在零点，则 \( f(x) \) 在区间 \( (1, \sqrt{\e}] \) 上仅有一个零点.
\end{enumerate}
\end{problem}

\begin{solution}
\begin{enumerate}
\item 由于 \(\ln x\) 的定义域为 \(x>0\)，函数 \(f\) 的定义域为 \((0,+\infty)\). 在此定义域内
\[
f'(x)=x-\frac{k}{x}=\frac{x^2-k}{x}
\]
因为 \(x>0\)，所以当 \(0<x<\sqrt{k}\) 时 \(f'(x)<0\)，当 \(x=\sqrt{k}\) 时 \(f'(x)=0\)，当 \(x>\sqrt{k}\) 时 \(f'(x)>0\). 因此 \(f(x)\) 在 \((0,\sqrt{k})\) 上单调递减，在 \((\sqrt{k},+\infty)\) 上单调递增，并在 \(x=\sqrt{k}\) 处取得极小值. 该极小值为
\[
f(\sqrt{k})=\frac{k}{2}-k\ln\sqrt{k}=\frac{k}{2}-\frac{k}{2}\ln k=\frac{k(1-\ln k)}{2}
\]
\item 若 \(f(x)\) 存在零点，则其最小值不大于 \(0\)，所以
\[
\frac{k(1-\ln k)}{2}\leq0
\]
由 \(k>0\) 得 \(1-\ln k\leq0\)，即 \(k\geq\e\).

当 \(k=\e\) 时，\(\sqrt{k}=\sqrt{\e}\)，函数在 \((1,\sqrt{\e})\) 上严格递减，且
\[
f(\sqrt{\e})=\frac{\e}{2}-\e\ln\sqrt{\e}=\frac{\e}{2}-\frac{\e}{2}=0
\]
所以此时区间 \((1,\sqrt{\e}]\) 内唯一的零点是 \(x=\sqrt{\e}\).

当 \(k>\e\) 时，\(\sqrt{\e}<\sqrt{k}\)，函数在 \([1,\sqrt{\e}]\) 上严格递减，并且
\(f(1)=\frac12>0\)，\(f(\sqrt{\e})=\frac{\e-k}{2}<0\).
由函数零点存在定理，\(f\) 在 \((1,\sqrt{\e})\) 内至少有一个零点；由严格单调性，该零点至多有一个. 因而在此情形下也恰有一个零点.

综上，若 \(f(x)\) 存在零点，则 \(f(x)\) 在区间 \((1,\sqrt{\e}]\) 上仅有一个零点.
\end{enumerate}
\end{solution}

\begin{problem}
已知椭圆 \(C: x^{2} + 3y^{2} = 3\)，过点 \(D(1,0)\) 且不过点 \(E(2,1)\) 的直线与椭圆 \(C\) 交于 \(A\)，\(B\) 两点，直线 \(AE\) 与直线 \(x = 3\) 交于点 \(M\).

\begin{enumerate}
\item 求椭圆 \(C\) 的离心率；
\item 若直线 \(AB\) 垂直于 \(x\) 轴，求直线 \(BM\) 的斜率；
\item 试判断直线 \(BM\) 与直线 \(DE\) 的位置关系，并说明理由.
\end{enumerate}
\end{problem}

\begin{solution}
\begin{enumerate}
\item 将椭圆方程化为标准形式，得
\[
\frac{x^2}{3}+y^2=1
\]
所以长半轴 \(a=\sqrt3\)，短半轴 \(b=1\)，焦半距 \(c=\sqrt{a^2-b^2}=\sqrt2\). 因此椭圆的离心率为
\[
e=\frac{c}{a}=\frac{\sqrt2}{\sqrt3}=\frac{\sqrt6}{3}
\]
\item 因为直线 \(AB\) 过 \(D(1,0)\) 且垂直于 \(x\) 轴，所以 \(AB\) 的方程为 \(x=1\). 设 \(A(1,y_1)\)，\(B(1,-y_1)\)，则由 \(E(2,1)\) 得直线 \(AE\) 的方程为
\[
y-1=(1-y_1)(x-2)
\]
令 \(x=3\)，得 \(M(3,2-y_1)\). 所以
\[
k_{BM}=\frac{(2-y_1)-(-y_1)}{3-1}=1
\]
\item 设直线 \(AB\) 为 \(\ell\). 先考虑 \(\ell\) 斜率不存在的情形，此时 \(\ell\) 为 \(x=1\)，由第（2）问知 \(k_{BM}=1\). 又
\[
k_{DE}=\frac{1-0}{2-1}=1
\]
所以 \(BM\parallel DE\).

下面考虑 \(\ell\) 斜率存在的情形. 设
\[
\ell:\ y=k(x-1)
\]
由于 \(\ell\) 不经过 \(E(2,1)\)，所以 \(k\ne1\). 设 \(A(x_1,y_1)\)，\(B(x_2,y_2)\)，则 \(y_i=k(x_i-1)\). 将直线方程代入椭圆方程，得
\[
(1+3k^2)x^2-6k^2x+3k^2-3=0
\]
由韦达定理，
\(x_1+x_2=\frac{6k^2}{1+3k^2}\)，\(x_1x_2=\frac{3k^2-3}{1+3k^2}\).
直线 \(AE\) 的方程为 \(y-1=\frac{y_1-1}{x_1-2}(x-2)\)，故令 \(x=3\) 得
\[
M\left(3,\frac{x_1+y_1-3}{x_1-2}\right)
\]
于是
\[
k_{BM}=\frac{\frac{x_1+y_1-3}{x_1-2}-y_2}{3-x_2}
\]
利用 \(y_1=k(x_1-1)\)、\(y_2=k(x_2-1)\)，化简得
\[
k_{BM}-1=\frac{(k-1)\left[-x_1x_2+2(x_1+x_2)-3\right]}{(3-x_2)(x_1-2)}
\]
再代入韦达定理，
\[
-x_1x_2+2(x_1+x_2)-3
=-\frac{3k^2-3}{1+3k^2}+\frac{12k^2}{1+3k^2}-3=0
\]
所以 \(k_{BM}=1=k_{DE}\)，从而 \(BM\parallel DE\). 综上，无论 \(AB\) 的斜率是否存在，直线 \(BM\) 均与直线 \(DE\) 平行.
\end{enumerate}
\end{solution}
